% GENERATED FILE. Edit canonical lessons, then run npm run generate:books.
% canonical-source-hash: fnv1a64:e28775ba61c14a26
% canonical-lessons: GU-C546-vagar, GU-C546-vina, GU-C546-sivay, GU-C546-uprat, GU-C546-kadi

\chapter{With or Without}
\label{ch:gu-with-or-without}
\hlchaptermodality{546}

\begin{chapteropening}
\textbf{By the end of this chapter:} \emph{I can say what happens with or without something.}

Complete the last lesson of chapter 546: I can say what happens with or without something.
\end{chapteropening}

\section[\texorpdfstring{\gu{વગર} (vagar)}{વગર (vagar)}]{\gu{વગર} (vagar) — without}

\label{lesson:GU-C546-vagar}



\begin{quote}
Before the new one: say the Gujarati for a try-out, then the Gujarati for a clever trick.
\end{quote}

\subsection*{You'll want to know: \gu{વગર}}
\textbf{\gu{વગર}} — \emph{vagar} — \textquotedblleft{}without\textquotedblright{}.

\begin{quote}
\gu{જો} \gu{તમે} \gu{છત્રી} \gu{વગર} \gu{જશો} \gu{તો} \gu{તમે} \gu{પલળી} \gu{જશો}.
\end{quote}

\begin{quote}
\emph{jo tame chatrī vagar jaśo to tame palḷī jaśo.}
\end{quote}

\begin{quote}
If you go without an umbrella, you'll get soaked.
\end{quote}

\begin{cousinweb}
\textbf{\gu{વગર}} (\emph{vagar}) is usually traced to Persian \emph{baghair}, \textquotedblleft{}without\textquotedblright{}. It follows its noun: \textbf{\gu{છત્રી} \gu{વગર}}.
\end{cousinweb}

\begin{grammarlens}[title={What you've built}]
One more word for saying what happens if.
\end{grammarlens}

\subsection*{Your turn}
\begin{itemize}
\raggedright
  \item \emph{Say it:} \emph{vagar}
  \item \emph{Say it:} \emph{vagar}, once more
  \item \emph{Say it:} \emph{jo tame chatrī vagar jaśo to tame palḷī jaśo.} — If you go without an umbrella, you'll get soaked.
  \item \emph{Recall:} say the Gujarati for a try-out, then the Gujarati for a clever trick, then say \emph{vagar} again
\end{itemize}

\begin{tcolorbox}[breakable,colback=teal!4,colframe=teal!35!black,title={Before you move on}]
What does \gu{વગર} mean? (\textquotedblleft{}Without\textquotedblright{}.) Say it once more.
\end{tcolorbox}
\section[\texorpdfstring{\gu{વિના} (vinā)}{વિના (vinā)}]{\gu{વિના} (vinā) — without (formal)}

\label{lesson:GU-C546-vina}



\begin{quote}
Before the new one: say the Gujarati for a clever trick, then the Gujarati for without.
\end{quote}

\subsection*{You'll want to know: \gu{વિના}}
\textbf{\gu{વિના}} — \emph{vinā} — \textquotedblleft{}without (formal)\textquotedblright{}.

\begin{quote}
\gu{જો} \gu{છોડ} \gu{પાણી} \gu{વિના} \gu{રહેશે} \gu{તો} \gu{તે} \gu{કરમાઈ} \gu{જશે}.
\end{quote}

\begin{quote}
\emph{jo choḍ pāṇī vinā raheśe to te karmāī jaśe.}
\end{quote}

\begin{quote}
If the plant goes without water, it will wilt.
\end{quote}

\begin{cousinweb}
\textbf{\gu{વિના}} (\emph{vinā}) goes back to Sanskrit \emph{vínā}, \textquotedblleft{}without\textquotedblright{} (Monier-Williams), which CDIAL lists with Gujarati forms (CDIAL 11772). It is a little more formal than \textbf{\gu{વગર}}.
\end{cousinweb}

\begin{grammarlens}[title={What you've built}]
One more word for saying what happens if.
\end{grammarlens}

\subsection*{Your turn}
\begin{itemize}
\raggedright
  \item \emph{Say it:} \emph{vinā}
  \item \emph{Say it:} \emph{vinā}, once more
  \item \emph{Say it:} \emph{jo choḍ pāṇī vinā raheśe to te karmāī jaśe.} — If the plant goes without water, it will wilt.
  \item \emph{Recall:} say the Gujarati for a clever trick, then the Gujarati for without, then say \emph{vinā} again
\end{itemize}

\begin{tcolorbox}[breakable,colback=teal!4,colframe=teal!35!black,title={Before you move on}]
What does \gu{વિના} mean? (\textquotedblleft{}Without (formal)\textquotedblright{}.) Say it once more.
\end{tcolorbox}
\section[\texorpdfstring{\gu{સિવાય} (sivāy)}{સિવાય (sivāy)}]{\gu{સિવાય} (sivāy) — except}

\label{lesson:GU-C546-sivay}



\begin{quote}
Before the new one: say the Gujarati for without, then the Gujarati for without (formal).
\end{quote}

\subsection*{You'll want to know: \gu{સિવાય}}
\textbf{\gu{સિવાય}} — \emph{sivāy} — \textquotedblleft{}except\textquotedblright{}.

\begin{quote}
\gu{જો} \gu{રવિવાર} \gu{સિવાય} \gu{કોઈ} \gu{દિવસ} \gu{તમને} \gu{અનુકૂળ} \gu{હશે} \gu{તો} \gu{મને} \gu{કહો}.
\end{quote}

\begin{quote}
\emph{jo ravivār sivāy koī divas tamne anukūḷ haśe to mane kaho.}
\end{quote}

\begin{quote}
If any day apart from Sunday suits you, tell me.
\end{quote}

\begin{cousinweb}
\textbf{\gu{સિવાય}} (\emph{sivāy}) is usually traced to Arabic \emph{siwā}, \textquotedblleft{}other than, except\textquotedblright{}, through Persian.
\end{cousinweb}

\begin{grammarlens}[title={What you've built}]
One more word for saying what happens if.
\end{grammarlens}

\subsection*{Your turn}
\begin{itemize}
\raggedright
  \item \emph{Say it:} \emph{sivāy}
  \item \emph{Say it:} \emph{sivāy}, once more
  \item \emph{Say it:} \emph{jo ravivār sivāy koī divas tamne anukūḷ haśe to mane kaho.} — If any day apart from Sunday suits you, tell me.
  \item \emph{Recall:} say the Gujarati for without, then the Gujarati for without (formal), then say \emph{sivāy} again
\end{itemize}

\begin{tcolorbox}[breakable,colback=teal!4,colframe=teal!35!black,title={Before you move on}]
What does \gu{સિવાય} mean? (\textquotedblleft{}Except\textquotedblright{}.) Say it once more.
\end{tcolorbox}
\section[\texorpdfstring{\gu{ઉપરાંત} (uprā̃t)}{ઉપરાંત (uprā̃t)}]{\gu{ઉપરાંત} (uprā̃t) — in addition to}

\label{lesson:GU-C546-uprat}



\begin{quote}
Before the new one: say the Gujarati for without (formal), then the Gujarati for except.
\end{quote}

\subsection*{You'll want to know: \gu{ઉપરાંત}}
\textbf{\gu{ઉપરાંત}} — \emph{uprā̃t} — \textquotedblleft{}in addition to\textquotedblright{}.

\begin{quote}
\gu{જો} \gu{પગાર} \gu{ઉપરાંત} \gu{ઘર} \gu{પણ} \gu{મળશે} \gu{તો} \gu{હું} \gu{આ} \gu{નોકરી} \gu{લઈશ}.
\end{quote}

\begin{quote}
\emph{jo pagār uprā̃t ghar paṇ maḷśe to hũ ā nokrī laīś.}
\end{quote}

\begin{quote}
If I get a house as well as the pay, I'll take this job.
\end{quote}

\begin{cousinweb}
\textbf{\gu{ઉપરાંત}} (\emph{uprā̃t}) follows its noun: \textbf{\gu{પગાર} \gu{ઉપરાંત}}, on top of the pay. No dictionary we use records an older form, so its history is not settled.
\end{cousinweb}

\begin{grammarlens}[title={What you've built}]
One more word for saying what happens if.
\end{grammarlens}

\subsection*{Your turn}
\begin{itemize}
\raggedright
  \item \emph{Say it:} \emph{uprā̃t}
  \item \emph{Say it:} \emph{uprā̃t}, once more
  \item \emph{Say it:} \emph{jo pagār uprā̃t ghar paṇ maḷśe to hũ ā nokrī laīś.} — If I get a house as well as the pay, I'll take this job.
  \item \emph{Recall:} say the Gujarati for without (formal), then the Gujarati for except, then say \emph{uprā̃t} again
\end{itemize}

\begin{tcolorbox}[breakable,colback=teal!4,colframe=teal!35!black,title={Before you move on}]
What does \gu{ઉપરાંત} mean? (\textquotedblleft{}In addition to\textquotedblright{}.) Say it once more.
\end{tcolorbox}
\section[\texorpdfstring{\gu{કદી} (kadī)}{કદી (kadī)}]{\gu{કદી} (kadī) — ever}

\label{lesson:GU-C546-kadi}



\begin{quote}
Before the new one: say the Gujarati for except, then the Gujarati for in addition to.
\end{quote}

\subsection*{You'll want to know: \gu{કદી}}
\textbf{\gu{કદી}} — \emph{kadī} — \textquotedblleft{}ever\textquotedblright{}.

\begin{quote}
\gu{જો} \gu{તમે} \gu{કદી} \gu{અમારા} \gu{ગામ} \gu{આવશો} \gu{તો} \gu{તમારે} \gu{અમારા} \gu{ઘરે} \gu{જ} \gu{રહેવું} \gu{પડશે}.
\end{quote}

\begin{quote}
\emph{jo tame kadī amārā gām āvśo to tamāre amārā ghare ja rahevũ paḍśe.}
\end{quote}

\begin{quote}
If you ever come to our village, you'll have to stay at our house.
\end{quote}

\begin{cousinweb}
\textbf{\gu{કદી}} (\emph{kadī}) is usually traced to Sanskrit \emph{kadā}, \textquotedblleft{}when?\textquotedblright{} (Monier-Williams). With a negative it is \textquotedblleft{}never\textquotedblright{}: \textbf{\gu{કદી} \gu{નહીં}}.
\end{cousinweb}

\begin{grammarlens}[title={What you've built}]
One more word for saying what happens if.
\end{grammarlens}

\subsection*{Your turn}
\begin{itemize}
\raggedright
  \item \emph{Say it:} \emph{kadī}
  \item \emph{Say it:} \emph{kadī}, once more
  \item \emph{Say it:} \emph{jo tame kadī amārā gām āvśo to tamāre amārā ghare ja rahevũ paḍśe.} — If you ever come to our village, you'll have to stay at our house.
  \item \emph{Recall:} say the Gujarati for except, then the Gujarati for in addition to, then say \emph{kadī} again
\end{itemize}

\begin{tcolorbox}[breakable,colback=teal!4,colframe=teal!35!black,title={Before you move on}]
What does \gu{કદી} mean? (\textquotedblleft{}Ever\textquotedblright{}.) Say it once more.
\end{tcolorbox}
